% META: {"id": "TEXP-MAT-CO-013", "chapitre": "TEXP-MATRICES-MARKOV", "type_objet": "corrige", "exercice_ref": "TEXP-MAT-EX-013", "capacites_codes": ["C6"], "sources_inspiration": [], "status": "generated"}
% BEGIN-VERIFY
% from sympy import Matrix, Rational, eye
% T = Matrix([[Rational(7, 10), Rational(3, 10)], [Rational(4, 10), Rational(6, 10)]])
% pi0 = Matrix([[Rational(1, 2), Rational(1, 2)]])
% pi = pi0
% for n in range(1, 9):
%     pi = pi * T
%     assert pi == pi0 * T ** n
% assert pi0 * T ** 0 == pi0 * eye(2) == pi0
% pi3 = pi0 * T ** 3
% assert pi3 == Matrix([[Rational(1139, 2000), Rational(861, 2000)]])
% END-VERIFY
\begin{corrige}{TEXP-MAT-EX-013}
\textbf{1.} La probabilité d'être en $j$ au rang $1$ s'obtient en sommant, sur
tous les états $i$ de départ, la probabilité d'être en $i$ multipliée par
celle de passer de $i$ à $j$ :
$(\pi_1)_j=\sum_i(\pi_0)_iT_{ij}$, ce qui est exactement le produit $\pi_0T$.

\textbf{2.} Au rang $0$ : $\pi_0=\pi_0T^0=\pi_0I$. Si $\pi_n=\pi_0T^n$, alors
en appliquant le même raisonnement qu'en $1$ entre les rangs $n$ et $n+1$ :
$\pi_{n+1}=\pi_nT=\pi_0T^nT=\pi_0T^{n+1}$. La propriété est héréditaire, donc
vraie pour tout $n$.

\textbf{3.} $\pi_3=\pi_0T^3=\begin{pmatrix}0{,}5695&0{,}4305\end{pmatrix}$, ce
que l'on retrouve en itérant trois fois $\pi\mapsto\pi T$.
\end{corrige}
